\documentclass[11pt,a4paper]{article}

\usepackage[T1]{fontenc}
\usepackage[utf8]{inputenc}
\usepackage[french]{babel}
\usepackage{lmodern}
\usepackage{geometry}
\geometry{margin=2.5cm}
\usepackage{amsmath,amssymb}
\usepackage{graphicx}
\usepackage{float}
\usepackage{hyperref}
\usepackage{caption}
\usepackage{adjustbox}

% Fichiers images/tables attendus à la racine du projet Overleaf

\title{Projet Monte Carlo - Pricing d'options et réduction de variance}
\author{Projet ST7}
\date{\today}

\begin{document}
\maketitle
\tableofcontents
\newpage
\section{Introduction}

La valorisation des options financières est un enjeu central en ingénierie financière. Si le modèle de Black-Scholes permet d'obtenir une formule fermée pour le prix d'une option européenne en temps continu, de nombreux produits dérivés ou situations de marché nécessitent des méthodes numériques. La méthode de Monte Carlo s'impose alors comme un outil incontournable pour estimer l'espérance de payoffs complexes, grâce à la simulation de scénarios aléatoires.

Cependant, la précision de l'estimation Monte Carlo dépend fortement de la variance de l'estimateur, et donc du nombre de simulations nécessaires pour obtenir un intervalle de confiance étroit. Pour améliorer l'efficacité de la méthode, diverses techniques de réduction de variance ont été développées. Ce projet a pour objectif d'étudier, d'implémenter et de comparer plusieurs de ces méthodes appliquées au pricing d'options européennes, d'abord en dimension 1, puis pour des baskets d'options.

Nous proposons également une étude de sensibilité des prix par rapport aux paramètres du modèle de Black-Scholes, afin d'analyser la robustesse des estimateurs et l'impact des paramètres sur la valorisation. L'ensemble des résultats est présenté sous forme d'une application interactive et de ce rapport technique, qui détaille les méthodes, les résultats numériques et leur interprétation.

\section{Contexte et objectifs}

La valorisation des options financières s'appuie sur des modèles mathématiques permettant d'estimer le prix d'un produit dérivé à partir de l'évolution stochastique de l'actif sous-jacent. Dans le cadre du cours ST7, plusieurs approches sont étudiées :

\begin{itemize}
    \item Le calcul stochastique en temps discret (arbre binomial), où le prix d'une option européenne peut être obtenu par récurrence, mais la formule à l'instant initial reste peu exploitable analytiquement.
    \item Le passage à la limite continue conduit au modèle de Black-Scholes, qui fournit une formule fermée pour le prix d'une option européenne, mais qui n'est pas généralisable à tous les produits.
    \item La méthode de Monte Carlo, fondée sur la simulation et le calcul d'espérances, permet d'estimer numériquement le prix d'options, même dans des situations complexes ou en grande dimension.
\end{itemize}

La méthode de Monte Carlo repose sur la loi forte des grands nombres et le théorème central limite (CLT), garantissant la convergence de l'estimateur et la possibilité de construire un intervalle de confiance. Cependant, la largeur de cet intervalle dépend de la variance de l'estimateur, d'où l'intérêt de techniques de réduction de variance.

L'objectif de ce projet est donc :
\begin{itemize}
    \item De comprendre et d'implémenter différentes méthodes de réduction de variance (antithétique, variable de contrôle, importance sampling -- implémenté ici via une translation gaussienne --, stratification, quasi-Monte Carlo, etc.) appliquées au pricing d'options européennes.
    \item De comparer leur efficacité sur des options de dimension 1 (où la formule de Black-Scholes sert de référence) puis sur des baskets d'options (dimension supérieure).
    \item De réaliser une étude de sensibilité du prix par rapport aux paramètres du modèle de Black-Scholes (spot, strike, volatilité, taux, maturité, corrélation, dimension).
    \item De présenter les résultats dans une application interactive et dans ce rapport technique, en mettant l'accent sur l'interprétation des résultats numériques.
\end{itemize}

Ce travail s'inscrit dans une démarche d'analyse critique des méthodes numériques utilisées en finance quantitative, avec une attention particulière portée à la robustesse, la précision et l'efficacité des estimateurs.

\section{Méthodes implémentées}

Dans ce projet, nous avons implémenté deux niveaux de méthodes :
\begin{itemize}
    \item un \textbf{socle commun} utilisé pour les comparaisons systématiques (notebook final + application Streamlit) ;
    \item des \textbf{extensions exploratoires} testées dans le notebook (optimisation plus poussée de l'importance sampling, variantes QMC supplémentaires).
\end{itemize}

Nous détaillons ci-dessous les méthodes effectivement codées et utilisées.

\subsection{Cadre de simulation commun}

Sous la mesure risque-neutre, nous simulons :
\begin{equation}
S_T = S_0\exp\left(\left(r-\frac{1}{2}\sigma^2\right)T + \sigma\sqrt{T}Z\right), \quad Z\sim\mathcal N(0,1).
\end{equation}
Le payoff actualisé d'un call européen est :
\begin{equation}
X = e^{-rT}(S_T-K)^+.
\end{equation}
Pour un panier arithmétique de dimension $d$, nous utilisons :
\begin{equation}
X = e^{-rT}\left(\sum_{i=1}^d w_i S_i(T)-K\right)^+,
\end{equation}
avec, dans les expériences principales, des poids uniformes $w_i=1/d$.

Toutes les méthodes rapportent le prix, l'erreur standard et l'IC95\%.

Dans la suite du document, nous utilisons les abréviations suivantes :
\begin{itemize}
    \item \textbf{MC} : Monte Carlo ;
    \item \textbf{RQMC} : Randomized Quasi-Monte Carlo ;
    \item \textbf{CV} : Control Variate (variable de contrôle) ;
    \item \textbf{IC95\%} : intervalle de confiance bilatéral à 95\% ;
    \item \textbf{PCA} : Principal Component Analysis (analyse en composantes principales).
\end{itemize}

\paragraph{Comment nous déterminons les incertitudes}
La quantité d'intérêt est la moyenne d'un estimateur de payoff actualisé. Pour un échantillon de taille $n$ (simulations MC ou réplications RQMC), nous estimons
\[
\widehat{\mathrm{SE}} = \frac{s}{\sqrt{n}},
\]
où $s$ est l'écart-type empirique de l'échantillon considéré.

L'intervalle de confiance à 95\% est construit par approximation gaussienne :
\[
\hat V \pm 1.96\,\widehat{\mathrm{SE}},
\]
et la largeur reportée dans les tableaux/figures vaut
\[
\mathrm{IC95\ width} = 2\times 1.96\times \widehat{\mathrm{SE}}.
\]

En pratique :
\begin{itemize}
    \item pour MC (standard, antithétique, CV, importance sampling par translation), l'échantillon est constitué des payoffs corrigés effectivement simulés ;
    \item pour RQMC, l'incertitude est estimée sur les moyennes de réplications randomisées indépendantes (scrambling), ce qui mesure la variabilité résiduelle de la méthode.
\end{itemize}

\subsection{Monte Carlo standard (MC)}

L'estimateur de base est :
\begin{equation}
\hat V_{MC} = \frac{1}{N}\sum_{i=1}^{N} X_i.
\end{equation}
L'IC95\% est construit par approximation gaussienne :
\begin{equation}
\hat V_{MC} \pm 1.96\,\frac{s}{\sqrt{N}},
\end{equation}
où $s$ est l'écart-type empirique des observations simulées.

\subsection{Variables antithétiques}

Nous exploitons la symétrie gaussienne via des paires $(Z,-Z)$, puis nous moyennons les payoffs par paires :
\begin{equation}
\hat V_{anti} = \frac{1}{N/2}\sum_{i=1}^{N/2}\frac{X(Z_i)+X(-Z_i)}{2}.
\end{equation}
Cette construction est présente en 1D et en basket, et constitue un composant clé de notre référence basket renforcée.

\subsection{Variable de contrôle (Control Variate)}

Nous corrigeons l'estimateur avec une variable de contrôle corrélée :
\begin{equation}
\hat V_{CV}=\bar X-\beta(\bar C-\mathrm{E}[C]),
\qquad
\beta^*=\frac{\mathrm{Cov}(X,C)}{\mathrm{Var}(C)}.
\end{equation}

En pratique :
\begin{itemize}
    \item en 1D, nous utilisons $C=e^{-rT}S_T$ avec $\mathrm{E}[C]=S_0$ ;
    \item en basket, nous utilisons $C=e^{-rT}B(T)$ avec $\mathrm{E}[C]=\sum_i w_iS_{0,i}$.
\end{itemize}

\subsection{Importance sampling par translation gaussienne}

\textbf{Clarification terminologique :} l'importance sampling est la méthode générale de changement de loi avec repondération (ratio de vraisemblance), tandis que la translation gaussienne est un cas particulier de cette méthode.

Dans l'application, nous implémentons précisément cette variante par translation gaussienne :
\begin{equation}
Z' = Z + \theta u,
\end{equation}
où $u$ est une direction normalisée (liée au gradient local du panier) et $\theta$ est sélectionné sur grille via un run pilote minimisant la variance empirique.

L'estimateur pondéré est :
\begin{equation}
\hat V_{IS}=\frac{1}{N}\sum_{i=1}^{N} X(Z_i+\theta u)\,\exp\left(-\theta u^\top Z_i-\frac{1}{2}\theta^2\|u\|^2\right).
\end{equation}

\subsection{Stratification}

Nous avons codé deux variantes :
\begin{itemize}
    \item \textbf{stratification naïve} sur la première composante gaussienne $Z_1$ ;
    \item \textbf{stratification PCA}, où la stratification se fait sur la projection principale de la structure de corrélation.
\end{itemize}
Ces deux versions sont implémentées dans l'application pour comparaison pédagogique et robustesse en basket.

\subsection{RQMC Sobol : variantes principales}

Nous utilisons des suites de Sobol randomisées (scrambling) avec $R$ réplications indépendantes afin d'estimer l'incertitude statistique sur la moyenne des réplications.

\paragraph{RQMC ICDF}
Chaque point $u\in(0,1)^d$ est transformé par $z=\Phi^{-1}(u)$ via l'\textbf{ICDF} (\textit{inverse cumulative distribution function}, c.-à-d. la fonction quantile de la loi normale), puis injecté dans la simulation Black-Scholes standard.

\paragraph{RQMC tronquée pondérée}
Chaque point est transformé par $z=-a+2au$ avec repondération par la densité gaussienne. En dimension $d$, les poids sont des produits de densités, ce qui explique une dégradation plus rapide en grande dimension (phénomène observé dans nos résultats).

\subsection{Génération des corrélations en basket}

Pour le cas multi-actifs, nous générons d'abord une matrice de corrélation cible (Identity, Constant, Toeplitz, Two-Block), puis :
\begin{equation}
Z^{corr}=ZL^\top, \quad \Sigma=LL^\top.
\end{equation}

Nous utilisons une robustification numérique explicite :
\begin{itemize}
    \item projection pseudo-\textbf{PSD} (\textit{positive semidefinite}, semi-définie positive) par clipping spectral des valeurs propres ;
    \item factorisation de Cholesky avec jitter diagonal progressif en cas de quasi-singularité.
\end{itemize}

\subsection{Méthodes exploratoires supplémentaires (notebook)}

En plus du socle principal, le notebook contient des extensions d'exploration :
\begin{itemize}
    \item importance sampling avec décalage optimisé (recherche numérique de $\mu$) ;
    \item importance sampling à deux paramètres $(\mu,\sigma_g)$ ;
    \item adaptation Cross-Entropy des paramètres de la loi d'échantillonnage ;
    \item variantes QMC basées sur Van der Corput et Box-Muller.
\end{itemize}
Ces briques ont servi à l'analyse méthodologique, mais nos comparaisons finales rapportées reposent principalement sur le socle commun (MC, antithétique, CV, importance sampling par translation, stratification, RQMC ICDF, RQMC tronquée pondérée), cohérent entre notebook final et application.

\section{Application interactive de pricing : conception, fonctionnalités et mode d'emploi}

Cette section décrit l'application interactive développée avec \texttt{Streamlit} (fichier \texttt{pricer\_app.py}). L'objectif est double :
\begin{itemize}
    \item fournir un environnement expérimental unifié pour comparer les méthodes Monte Carlo et RQMC ;
    \item permettre une exploration rapide de l'impact des paramètres de marché et de structure (dimension, corrélation, budget de simulation) sur le prix et la précision statistique.
\end{itemize}

L'application est accessible ici :
\href{https://moncomblehugo-monte-carlo-project-s8-pricer-app-nto6fh.streamlit.app/}{https://moncomblehugo-monte-carlo-project-s8-pricer-app-nto6fh.streamlit.app/}.

L'application reproduit les méthodes étudiées dans le notebook dans un cadre plus opérationnel : choix du produit, paramétrage, lancement des simulations, puis affichage direct de tableaux comparatifs, intervalles de confiance et visualisations.

\subsection{Organisation générale de l'interface}

Nous avons structuré l'application en deux blocs principaux :
\begin{itemize}
    \item une barre latérale de configuration ;
    \item deux onglets d'analyse (\textit{Pricing} et \textit{Sensibilité}).
\end{itemize}

\paragraph{Barre latérale (inputs)}
Dans la barre latérale, nous centralisons tous les paramètres de calcul :
\begin{itemize}
    \item \textbf{type de produit} : option européenne 1D ou panier arithmétique ;
    \item \textbf{paramètres financiers} : $S_0$, $K$, $T$, $r$, $\sigma$ ;
    \item \textbf{dimension} $d$ du panier (jusqu'à plusieurs dizaines d'actifs) ;
    \item \textbf{structure de corrélation} ;
    \item \textbf{budget de calcul} : $N_{MC}$, $N_{RQMC}$, nombre de réplications RQMC ;
    \item \textbf{sélection des méthodes} à comparer ;
    \item \textbf{seed} pour la reproductibilité.
\end{itemize}

\paragraph{Onglet Pricing}
Dans cet onglet, nous affichons :
\begin{itemize}
    \item la matrice de corrélation utilisée (carte de chaleur) ;
    \item un tableau comparatif (prix, erreur standard, IC95\%, temps de calcul) ;
    \item des graphiques de comparaison (prix, largeur d'IC95\%, écart à une référence).
\end{itemize}

\paragraph{Onglet Sensibilité}
Dans cet onglet, nous faisons varier un paramètre (ex : $S_0$, $K$, $\sigma$, $r$, $T$, $d$, $\rho$) sur une grille et nous observons :
\begin{itemize}
    \item l'évolution du prix ;
    \item l'évolution de la largeur d'intervalle de confiance ;
    \item la stabilité des méthodes selon les zones de paramètres.
\end{itemize}

\subsection{Construction des scénarios corrélés}

La logique de corrélation reprend le cadre méthodologique présenté en section précédente ; l'accent est mis ici sur les choix d'implémentation utiles à l'usage interactif.

\paragraph{Étape 1 : tirages indépendants}
Nous générons d'abord des vecteurs $Z \sim \mathcal N(0, I_d)$.

\paragraph{Étape 2 : matrice de corrélation cible}
Nous définissons une matrice $\Sigma$ selon la structure choisie (voir section suivante).

\paragraph{Étape 3 : factorisation de Cholesky}
Nous calculons une factorisation $L$ telle que :
\begin{equation}
\Sigma = LL^\top,
\end{equation}
puis nous construisons les facteurs corrélés :
\begin{equation}
Z^{corr} = ZL^\top.
\end{equation}

\paragraph{Étape 4 : robustification numérique}
En pratique, certaines matrices proches de la singularité peuvent rendre la factorisation instable. Le traitement appliqué est le suivant :
\begin{itemize}
    \item une \textbf{projection pseudo-PSD} (clipping spectral des valeurs propres négatives) ;
    \item un \textbf{jitter diagonal} progressif ($\epsilon I$) en cas d'échec de Cholesky.
\end{itemize}
Cette combinaison évite les plantages sur des configurations extrêmes (fortes corrélations, grandes dimensions).

\subsection{Matrices de corrélation proposées dans l'application}

Nous proposons plusieurs structures de dépendance :

\paragraph{1) Identity}
\begin{equation}
\Sigma = I_d.
\end{equation}
Les actifs sont indépendants conditionnellement aux paramètres individuels.

\paragraph{2) Constant}
Toutes les corrélations hors diagonale sont égales à $\rho$ :
\begin{equation}
\Sigma_{ij}=1 \quad \mbox{si } i=j, \qquad
\Sigma_{ij}=\rho \quad \mbox{si } i\neq j.
\end{equation}
Nous appliquons un clipping pour respecter la condition de validité approximative $\rho > -1/(d-1)$.

\paragraph{3) Toeplitz}
La corrélation décroît avec la distance entre indices :
\begin{equation}
\Sigma_{ij} = \rho^{|i-j|}.
\end{equation}
Cette structure modélise des actifs ordonnés (maturités, secteurs proches, etc.).

\paragraph{4) Two-Block}
Les actifs sont séparés en deux blocs, avec corrélation intra-bloc et inter-blocs :
\begin{itemize}
    \item corrélation \textit{intra} contrôlée par $\rho$ ;
    \item corrélation \textit{inter} contrôlée par $\rho_2$.
\end{itemize}
Ce modèle permet de représenter des univers multi-secteurs ou multi-régimes.

\subsection{Gestion des options panier}

Pour un panier arithmétique de dimension $d$, nous simulons les valeurs terminales
\begin{equation}
S_i(T) = S_{0,i}\exp\left(\left(r - \frac{1}{2}\sigma_i^2\right)T + \sigma_i\sqrt{T}\,Z_i^{corr}\right),
\end{equation}
et nous construisons le panier
\begin{equation}
B(T) = \sum_{i=1}^d w_i S_i(T),
\end{equation}
avec, dans la version courante, des poids uniformes $w_i = 1/d$.

Le payoff actualisé est alors :
\begin{equation}
X = e^{-rT}(B(T)-K)^+.
\end{equation}

Ce cadre permet de faire varier facilement la dimension et la corrélation pour étudier les effets de \textit{curse of dimensionality} sur la variance et la stabilité numérique.

\subsection{Méthodes disponibles dans l'application}

Nous intégrons un ensemble cohérent de méthodes, toutes évaluées avec les mêmes entrées pour une comparaison équitable.

\paragraph{MC standard}
Nous utilisons cet estimateur de base comme référence opérationnelle.

\paragraph{MC antithétique}
Nous utilisons les paires $(Z,-Z)$ pour réduire la variance via corrélation négative.

\paragraph{MC + variable de contrôle}
Nous utilisons comme contrôle la valeur actualisée du panier, dont l'espérance est connue :
\begin{equation}
\mathrm{E}\left[e^{-rT}B(T)\right] = \sum_i w_i S_{0,i}.
\end{equation}

\paragraph{MC antithétique + variable de contrôle}
Cette version est généralement plus performante, et nous l'utilisons comme référence interne en basket quand Black-Scholes n'est pas disponible.

\paragraph{MC translation (variante d'importance sampling gaussien)}
Nous décalons les tirages dans une direction pertinente (dérivée du gradient du panier) ; le paramètre de translation est choisi sur grille pilote pour minimiser la variance empirique.

\paragraph{RQMC Sobol + ICDF}
Nous transformons les points Sobol par $\Phi^{-1}$ ; cette version est robuste en dimension modérée à élevée et fournit souvent les IC95\% les plus serrés à budget donné.

\paragraph{RQMC tronquée pondérée}
Nous effectuons une transformation uniforme vers $[-a,a]^d$ puis une repondération par la densité normale. Cette méthode est efficace en faible dimension, mais plus fragile lorsque $d$ augmente (poids produits très dispersés).

\paragraph{Stratification naïve et PCA}
Nous proposons deux variantes de stratification pour la comparaison pédagogique :
\begin{itemize}
    \item stratification sur la première composante ;
    \item stratification sur la composante principale (PCA) de la corrélation.
\end{itemize}

\subsection{Mesures reportées et logique de comparaison}

Pour chaque méthode, nous reportons systématiquement :
\begin{itemize}
    \item le prix estimé ;
    \item l'erreur standard ;
    \item l'intervalle de confiance à 95\% ;
    \item la largeur d'IC95\% (indicateur principal de précision) ;
    \item le temps d'exécution.
\end{itemize}

Lorsque c'est possible, une référence est définie :
\begin{itemize}
    \item en dimension 1 : prix exact Black-Scholes ;
    \item en basket : estimateur MC antithétique + variable de contrôle exécuté avec un budget élevé, retenu comme meilleure approximation interne.
\end{itemize}
Nous traçons ensuite l'écart absolu à cette référence pour faciliter l'analyse biais/variance.

\section{Résultats numériques complets et analyse détaillée}

Les résultats de cette section proviennent du bloc final du notebook (graphes et tableaux générés automatiquement dans \texttt{figures/}), complétés par une interprétation quantitative.

\subsection{Cadre d'évaluation}

Nous évaluons les méthodes principales sur trois axes :
\begin{itemize}
    \item \textbf{précision statistique} : largeur d'IC95\% ;
    \item \textbf{justesse} : erreur absolue (vs Black-Scholes en 1D, vs référence MC renforcée en basket) ;
    \item \textbf{coût} : temps de calcul.
\end{itemize}

Le benchmark est réalisé sur plusieurs dimensions : $d\in\{1,2,5,10,20,40\}$.

En dimension strictement supérieure à 1 (cas panier), il n'existe pas de formule fermée de référence de type Black-Scholes. La référence numérique retenue est donc un \textbf{MC renforcé} (antithétique + variable de contrôle) exécuté avec un \textbf{très grand nombre de simulations}. Ce choix est motivé par un compromis favorable entre coût et réduction de variance : la méthode reste suffisamment peu coûteuse pour être lancée sur un budget élevé, ce qui fournit une approximation très précise du prix de référence. À budget équivalent, en dimension 1, cette approche atteint typiquement une erreur de l'ordre de $10^{-7}$, ce qui conforte son utilisation comme proxy de référence en dimension supérieure.

Dans toute la section, \textbf{constantes initiales} désigne le jeu de paramètres de référence utilisé quand un paramètre n'est pas en cours de variation :
\[
S_0=100,\quad K=100,\quad \sigma=0.20,\quad r=0.10,\quad T=1.0.
\]
Pour les tests panier, on utilise des poids uniformes $w_i=1/d$, une corrélation constante de niveau $\rho=0.30$ et la même volatilité par actif (cas homogène), sauf mention contraire.

\subsection{Sensibilité des méthodes aux paramètres}

Pour chaque paramètre, les courbes \textbf{prix}, \textbf{IC95\%} et \textbf{écart absolu à la référence Black-Scholes} sont tracées pour les méthodes principales (MC, MC+CV, RQMC ICDF, RQMC tronquée).

\textbf{Important :} les cinq graphes de sensibilité ci-dessous sont réalisés en \textbf{dimension 1} (option européenne 1D), en faisant varier un seul paramètre à la fois autour des constantes de base.

\begin{figure}[h!]
    \centering
    \includegraphics[width=0.9\textwidth]{results_refait_sens_S0.png}
    \caption{Sensibilité des méthodes à $S_0$ (prix, IC95\% et écart absolu à la référence).}
    \label{fig:results_refait_sens_s0}
\end{figure}

\begin{figure}[h!]
    \centering
    \includegraphics[width=0.9\textwidth]{results_refait_sens_K.png}
    \caption{Sensibilité des méthodes à $K$ (prix, IC95\% et écart absolu à la référence).}
    \label{fig:results_refait_sens_k}
\end{figure}

\begin{figure}[h!]
    \centering
    \includegraphics[width=0.9\textwidth]{results_refait_sens_sigma.png}
    \caption{Sensibilité des méthodes à $\sigma$ (prix, IC95\% et écart absolu à la référence).}
    \label{fig:results_refait_sens_sigma}
\end{figure}

\begin{figure}[h!]
    \centering
    \includegraphics[width=0.9\textwidth]{results_refait_sens_r.png}
    \caption{Sensibilité des méthodes à $r$ (prix, IC95\% et écart absolu à la référence).}
    \label{fig:results_refait_sens_r}
\end{figure}

\begin{figure}[h!]
    \centering
    \includegraphics[width=0.9\textwidth]{results_refait_sens_T.png}
    \caption{Sensibilité des méthodes à $T$ (prix, IC95\% et écart absolu à la référence).}
    \label{fig:results_refait_sens_t}
\end{figure}

Lecture quantitative principale :
\begin{itemize}
    \item les sensibilités économiques attendues sont respectées (prix croissant en $S_0$, $\sigma$, $r$, $T$ et décroissant en $K$) ;
    \item les méthodes RQMC donnent des IC95\% nettement plus serrés que MC sur la plupart des plages ;
    \item la variante tronquée peut être excellente en 1D, mais devient instable dès que la dimension augmente (confirmé dans le benchmark multi-dimension ci-dessous).
\end{itemize}

\subsection{Comparaison précision / justesse / runtime selon la dimension}

Le tableau complet des métriques est donné ci-dessous.

Les trois graphes suivants sont des graphes \textbf{multi-dimensions} : l'axe horizontal correspond à $d\in\{1,2,5,10,20,40\}$. Le point $d=1$ correspond au cas 1D (référence Black-Scholes), tandis que les points $d\ge 2$ correspondent au cas panier (référence MC renforcée).

\begin{table}[h!]
    \centering
    \caption{Métriques complètes par méthode et dimension (prix, IC95, runtime, erreur absolue).}
    \label{tab:results_refait_full}
    \begin{adjustbox}{max width=\textwidth}
    \input{results_refait_table_full.tex}
    \end{adjustbox}
\end{table}

\begin{figure}[h!]
    \centering
    \includegraphics[width=0.78\textwidth]{results_refait_ci_vs_dim.png}
    \caption{Largeur d'IC95\% selon la dimension.}
    \label{fig:results_refait_ci_vs_dim}
\end{figure}

\begin{figure}[h!]
    \centering
    \includegraphics[width=0.78\textwidth]{results_refait_abs_error_vs_dim.png}
    \caption{Erreur absolue selon la dimension.}
    \label{fig:results_refait_abs_error_vs_dim}
\end{figure}

\begin{figure}[h!]
    \centering
    \includegraphics[width=0.78\textwidth]{results_refait_runtime_vs_dim.png}
    \caption{Temps de calcul selon la dimension.}
    \label{fig:results_refait_runtime_vs_dim}
\end{figure}

Résultats numériques marquants :
\begin{itemize}
    \item en \textbf{1D}, RQMC tronquée est la plus précise et la plus juste : IC95\% $=9.53\times 10^{-6}$, erreur absolue $=2.06\times 10^{-6}$ ;
    \item en \textbf{1D}, MC+CV est la plus rapide parmi les méthodes comparées (runtime $\approx 2.91\times 10^{-3}$ s), avec une bonne justesse (erreur $\approx 1.27\times 10^{-3}$) ;
    \item en \textbf{dimension intermédiaire} ($d=5,10$), MC reste compétitif en compromis global (score synthétique minimal), malgré une précision statistique inférieure à RQMC ;
    \item en \textbf{grande dimension} ($d=20,40$), RQMC ICDF devient la méthode la plus robuste (IC95\% de l'ordre de $10^{-2}$ et erreur absolue faible), tandis que la variante tronquée s'effondre (prix proche de 0 pour $d\ge 20$ dans nos tests).
\end{itemize}

\subsection{Tableau de synthèse des meilleures méthodes}

Nous condensons enfin les résultats via un score compromis (justesse + précision + coût, normalisé par dimension).

\begin{table}[h!]
    \centering
    \caption{Meilleure méthode par dimension (score compromis).}
    \label{tab:results_refait_best}
    \begin{adjustbox}{max width=\textwidth}
    \input{results_refait_best_by_dim.tex}
    \end{adjustbox}
\end{table}

Lecture rapide :
\begin{itemize}
    \item $d=1$ : RQMC ICDF ;
    \item $d=2$ : RQMC tronquée ;
    \item $d=5,10$ : MC ;
    \item $d=20,40$ : RQMC ICDF.
\end{itemize}

Conclusion opérationnelle de la section résultats :
\begin{itemize}
    \item pour l'option 1D, RQMC tronquée domine en précision/justesse, tandis que RQMC ICDF ressort meilleur en score compromis ;
    \item pour les dimensions élevées, RQMC ICDF est la méthode la plus fiable ;
    \item pour des dimensions intermédiaires avec contrainte temps, MC peut rester le meilleur compromis pratique.
\end{itemize}

\subsection{Comparaison finale poussée : dimension 1 vs dimension 30}

En complément du benchmark global, nous ajoutons deux tableaux comparatifs finaux construits sur un grand nombre de points avec les constantes initiales :
\begin{itemize}
    \item un tableau en \textbf{dimension 1} (avec référence Black-Scholes) ;
    \item un tableau en \textbf{dimension 30} (cas panier, avec référence interne MC antithétique + CV à grand budget).
\end{itemize}

Ces tableaux permettent une lecture directe des compromis précision/temps dans deux régimes extrêmes : faible dimension et dimension élevée.

\begin{figure}[h!]
    \centering
    \includegraphics[width=0.95\textwidth]{tableau_comparatif_dim1.png}
    \caption{Tableau comparatif final en dimension 1 (constantes initiales).}
    \label{tab:comparatif_final_dim1}
\end{figure}

\begin{figure}[h!]
    \centering
    \includegraphics[width=0.95\textwidth]{tableau_comparatif_dim30.png}
    \caption{Tableau comparatif final en dimension 30 (constantes initiales).}
    \label{tab:comparatif_final_dim30}
\end{figure}

Lecture synthétique :
\begin{itemize}
    \item en dimension 1, les méthodes RQMC obtiennent les meilleures précisions absolues, au prix d'un temps de calcul plus élevé ;
    \item en dimension 30, RQMC ICDF reste précis mais coûteux, et la variante tronquée devient instable (prix quasi nul), ce qui confirme les observations du benchmark multi-dimension ;
    \item MC Anti + CV reste une référence robuste en grande dimension pour maintenir un bon compromis opérationnel.
\end{itemize}

\subsection{Validation du sélecteur automatique de méthode}

L'objectif de cette partie est d'évaluer une stratégie de sélection automatique de méthode de pricing, adaptée à un usage interactif : obtenir rapidement une estimation fiable sans exécuter systématiquement l'ensemble des pricers disponibles.

\paragraph{1) Comment le sélecteur Auto décide dans l'application}
La logique implémentée dans \texttt{pricer\_app.py} (fonction \texttt{select\_best\_method\_auto}) repose sur un arbitrage explicite \textbf{précision vs temps de calcul}.

Le schéma de décision est le suivant :
\begin{enumerate}
    \item définir un sous-ensemble de méthodes candidates en fonction du régime de dimension ;
    \item exécuter un \textbf{mini-benchmark pilote} avec des budgets réduits ;
    \item mesurer, pour chaque candidate, la largeur d'IC95\% (précision) et le runtime (coût) ;
    \item normaliser ces deux métriques et calculer un score unique :
    \[
    \mathrm{Score}=\omega\,\widetilde{\mathrm{IC95}}+(1-\omega)\,\widetilde{\mathrm{Runtime}},
    \]
    où $\omega\in[0,1]$ est un poids réglable (0 : priorité vitesse, 1 : priorité précision) ;
    \item sélectionner la méthode de score minimal.
\end{enumerate}

Cette formulation rend le compromis opérationnel explicite et facilement interprétable.

Concrètement, dans l'application, le déroulé est le suivant :
\begin{enumerate}
    \item l'utilisateur choisit le mode \textit{Auto (choix optimal)} et règle le curseur de poids de précision $\omega$ ;
    \item lors du clic sur \textit{Lancer le pricing}, l'app construit l'ensemble des candidates autorisées en mode Auto (MC, MC Translation, MC Anti+CV, RQMC ICDF, RQMC Truncated, Stratif. PCA) ;
    \item cet ensemble est filtré selon la dimension :
    \begin{itemize}
        \item $d=1$ : \{MC Anti+CV, MC Translation, RQMC ICDF, RQMC Truncated, MC+CV\} ;
        \item $2\le d\le 20$ : \{MC Anti+CV, RQMC ICDF, Stratif. PCA, MC\} ;
        \item $d>20$ : \{MC Anti+CV, RQMC ICDF, MC\}.
    \end{itemize}
    Si l'intersection est vide, l'app force un fallback sécurisé vers MC Anti+CV ;
    \item l'app exécute un mini-benchmark pilote avec des budgets réduits, calculés automatiquement à partir des budgets demandés :
    \[
    N_{MC}^{pilot}=\min\!\left(30000,\max\!\left(6000,\left\lfloor\frac{N_{MC}}{20}\right\rfloor\right)\right),
    \]
    \[
    N_{RQMC}^{pilot}=\min\!\left(4096,\max\!\left(1024,\left\lfloor\frac{N_{RQMC}}{8}\right\rfloor\right)\right),
    \]
    \[
    R^{pilot}=\min\!\left(8,\max\!\left(4,\left\lfloor\frac{R}{2}\right\rfloor\right)\right);
    \]
    \item pour chaque méthode candidate, l'app mesure la largeur d'IC95\% et le runtime pilote, puis normalise chaque métrique par min-max sur les candidates ;
    \item le score final est calculé par
    \[
    \mathrm{Score}=\omega\,\widetilde{\mathrm{IC95}}+(1-\omega)\,\widetilde{\mathrm{Runtime}},
    \]
    et la méthode de score minimal est retenue ;
    \item la méthode choisie est ensuite exécutée avec le budget complet demandé par l'utilisateur pour produire le prix final.
\end{enumerate}

Le tableau de diagnostic affiché dans l'application reprend précisément ces éléments (rang, IC95\% pilote, runtime pilote, composantes normalisées, score et méthode sélectionnée), ce qui rend la décision traçable pour l'utilisateur.

\paragraph{Paramètres de calibration}
Le calibrage porte sur :
\begin{itemize}
    \item les budgets du mini-benchmark pilote ;
    \item la liste des méthodes candidates par régime de dimension ;
    \item le poids $\omega$ du compromis précision/temps.
\end{itemize}

Le choix d'un mécanisme court de benchmarking pilote limite le risque de sur-ajustement à des règles complexes et reste robuste aux variations d'environnement de calcul.

Enfin, un \textbf{fallback de sécurité} conserve MC Anti + CV si une méthode hors liste autorisée apparaît.
Le diagnostic affiché dans l'application fournit les métriques pilotes (IC95\%, runtime), les scores normalisés et la méthode retenue.

\paragraph{2) Ce que valide le benchmark de fin de rapport}
La validation est effectuée sur le scénario de référence (constantes initiales), avec un compromis équilibré $\omega=0.5$ entre précision et temps de calcul. Le tableau ci-dessous (issu de \texttt{automode\_verif.csv}) fournit les métriques finales utiles à la décision : prix estimé, précision statistique (SE, IC95\%), coût de calcul et écart à la référence.

\begin{table}[h!]
    \centering
    \caption{Validation du mode automatique (constantes initiales, $\omega=0.5$).}
    \label{tab:automode_verif}
    \begin{adjustbox}{max width=\textwidth}
    \input{automode_verif.tex}
    \end{adjustbox}
\end{table}

Cette validation permet de vérifier simultanément deux objectifs :
\begin{itemize}
    \item \textbf{qualité de décision} : la méthode effectivement choisie conserve une précision compatible avec l'objectif de l'utilisateur ;
    \item \textbf{efficacité temporelle} : le runtime reste cohérent avec un usage interactif.
\end{itemize}

Autrement dit, nous déplaçons l'effort coûteux vers l'analyse offline et nous conservons en ligne une décision rapide, traçable et robuste. Dans certains scénarios, cette accélération peut impliquer un compromis mesuré sur la précision : c'est un choix assumé pour un usage interactif, avec possibilité de basculer en mode manuel quand l'utilisateur vise la meilleure précision possible.

Nous obtenons ainsi un outil de pricing réellement opérationnel : l'utilisateur fournit les paramètres du produit, et l'application choisit automatiquement une méthode pertinente avant de produire le prix final avec ses indicateurs de confiance.


\section{Conclusion}

Ce projet a permis de construire un cadre complet de valorisation d'options par simulation, allant de l'implémentation des méthodes de base jusqu'à leur comparaison systématique dans un environnement reproductible (notebook) et opérationnel (application Streamlit).

Sur le plan méthodologique, les résultats confirment les enseignements attendus :
\begin{itemize}
    \item en dimension 1, les approches RQMC atteignent les meilleures précisions statistiques à budget donné ;
    \item en dimension intermédiaire, le compromis précision/temps peut rester favorable à des variantes MC bien contrôlées ;
    \item en grande dimension, RQMC ICDF apparaît comme l'option la plus robuste parmi les méthodes testées ;
    \item la combinaison antithétique + variable de contrôle constitue une référence numérique fiable en l'absence de formule fermée.
\end{itemize}

Au-delà des performances ponctuelles, la contribution principale du travail est d'avoir relié trois niveaux complémentaires :
\begin{itemize}
    \item \textbf{analyse quantitative} (précision, justesse, runtime) ;
    \item \textbf{interprétation financière} (sensibilités aux paramètres, effets de dimension et de corrélation) ;
    \item \textbf{décision opérationnelle} (sélecteur automatique de méthode, explicable et rapide).
\end{itemize}

Le projet présente néanmoins des limites naturelles : calibration empirique des seuils du sélecteur, dépendance au budget machine, et périmètre produit centré sur des payoffs européens/paniers arithmétiques. Ces limites n'affectent pas la cohérence des conclusions, mais indiquent les points à renforcer pour une généralisation.

Les perspectives les plus directes sont les suivantes :
\begin{itemize}
    \item étendre l'étude à d'autres payoffs (barrières, asiatiques, max/min) ;
    \item formaliser un recalibrage automatique des règles de sélection selon l'environnement de calcul ;
    \item intégrer des diagnostics d'erreur en ligne plus fins (stabilité des IC95\%, détection de scénarios difficiles) ;
    \item comparer avec des schémas de simulation supplémentaires (multilevel Monte Carlo, variantes QMC avancées).
\end{itemize}

En synthèse, le rapport et l'application démontrent qu'un pipeline Monte Carlo bien structuré peut fournir des prix fiables, des indicateurs d'incertitude exploitables et une logique de choix de méthode robuste, ce qui répond à l'objectif initial d'un outil de pricing à la fois pédagogique et utilisable en pratique.








\end{document}

